% 
% ###################################################################
% Copyright (c) 2018, AuthorName 
%  Editors: A.D. Rollett & M. De Graef
% All rights reserved.
%
% Licensed under the Creative Commons CC BY-NC-SA 4.0 License, 
% hereafter referred to as the "License"; you may not use this 
% document except in compliance with the License. You may obtain 
% a copy of the License at 
%     https://creativecommons.org/licenses/by-nc-sa/4.0/legalcode 
% Unless required by applicable law or agreed to in writing, all 
% material distributed under the License is distributed on an 
% "AS IS" BASIS, WITHOUT WARRANTIES OR CONDITIONS OF ANY KIND, 
% either express or implied. See the License for the specific 
% language governing permissions and limitations under the License.
% ###################################################################

% ###################################################################
% The following lines are to be uncommented or edited as needed 

%\corechapter{Yes}   % uncomment this line only if this chapter is a core/foundational chapter;
\OCchapterauthor{Marc De Graef, Carnegie Mellon University} %this will appear in a secondary header below the chapter title.

% graphics path for all figure *.pdf files 

%	redefine the \chabbr command to be a six-character string that uniquely identifies this 
%	chapter; easiest way to do this is to take the first couple of letters from the title words
\renewcommand{\chabbr}{ORISAM}

%     replace \noheaderimage by the chapter header image file name (without .pdf extension).
%     Chapter header images must be 2480 x 1240 pixels with 300dpi, RGB format.
\chapterimage{ORISAMheader.pdf}

%	Make sure to use the \label{\chabbr:} command everywhere instead of \label{} !
\chapter{Orientation Sampling and Averaging}\label{\chabbr:ch:OrSamp}

% add the author information so that it will appear in the Author List at the start of the document
\writeauthor{ORISAM:ch:OrSamp}{Orientation Sampling and Averaging}{De Graef}{Marc}{Materials Science and Engineering}{Carnegie Mellon University}{mdg@andrew.cmu.edu}{https://www.mse.engineering.cmu.edu/directory/bios/degraef-marc.html}

% Each chapter begins with Learning Objectives; the list of objectives should have links to sections/subsections using their OC labels
% each Learning Objective should be an active statement (i.e., contain a verb).  The \\ command can be used to force an item into the 
% second column if LaTeX breaks the line at an awkward location.
\begin{learningobjectives}
{\begin{itemize}
    \item[{\color{OCBurntOrange}\ref{\chabbr:sec:sampling}:}] Learn how sample rotations/orientations
    \item[{\color{OCBurntOrange}\ref{\chabbr:sec:Riesz}:}] Learn how to evaluate sampling uniformity
    \item[{\color{OCBurntOrange}\ref{\chabbr:sec:orientationaveraging}:}] Describe averaging in periodic spaces 
    \item[{\color{OCBurntOrange}\ref{\chabbr:sec:samplingRFZs}:}] Sample and average orientations in a fundamental zone 
\end{itemize}}
\end{learningobjectives}

% ###################################################################
% ###################################################################
% ###################################################################

\section{Introductory Remarks}\label{\chabbr:sec:intro}



\section{Sampling of Orientations/Rotations}\label{\chabbr:sec:sampling}

In this section, we take a closer look at different ways to generate random sets of orientations that are either uniform or clustered. The techniques described here belong to the field of \indexit{Directional Statistics} which handles the theory of uniform and clustered distributions of points on (hyper-)spheres.  While the following sections present a basic overview, the interested reader is encouraged to consult the excellent book ``Directional Statistics'' by Mardia and Jupp \cite{mardia2009a}. Before we cover these sampling concepts, it is important to introduce the concept of \indexit{misorientation}, which is a measure for the ``distance'' between two orientations.

\subsection{Misorientations}\label{\chabbr:ssec:misorientations}

We have defined a rotation as a passive operation that rotates one cartesian reference frame (the ``old'' frame) into another reference frame (the ``new'' frame).  In terms of rotation matrices, we commonly write:
\[
	\mathbf{e}'_i = \alpha_{ij}\mathbf{e}_j\ ,
\]
where $\alpha_{ij}\equiv \mathbf{e}'_i\cdot\mathbf{e}_j$ is a special orthogonal $3\times 3$ matrix, i.e., it belongs to the rotation group $SO(3)$. It is common practice in the texture community to select a reference frame attached to the sample as the old reference frame, and the reference frame of an individual grain as the new reference frame.  In addition to the rotation matrix $\alpha_{ij}$ we can use Euler angles or any of the neo-Eulerian representations introduced before to describe the grain orientation.  In this subsection, we will use the quaternion representation from here on out.

A polycrystalline material has many grains that form pairwise \indexit{grain boundaries}; each grain orientation can be represented by a quaternion $q_i$ with $i\in[1,N]$ with $N$ the total number of grains.  If we focus on two grains that are nearest neighbors, say grains $1$ and $2$, then their orientations with respect to the sample reference frame are given by the quaternions $q_1$ and $q_2$.  We now ask the question: how can we rotate directly from the reference frame of grain $1$ to that of grain $2$?  To answer this question we can make use of the fact that the inverse of a rotation is described by the complex conjugate of the rotation quaternion.  So, to get from grain $1$ to grain $2$ (in terms of their orientation), we can first go from grain $1$ to the sample reference frame by means of $q_1^{\ast}$, and then from the sample frame to grain $2$ by means of $q_2$.  This results in the \indexit{misorientation quaternion} $q_m$ (ordering the passive quaternions from left to right):
\begin{equation}
	q_m = q_1^{\ast}\,q_2\rightarrow q_m^{\ast} = (q_1^{\ast}\,q_2)^{\ast} = q_2^{\ast}\,q_1\ .\label{\chabbr:eq:misorientation}
\end{equation}
The misorientation quaternion $q_m$ rotates the reference frame of grain $1$ into the orientation of the reference frame of grain $2$.  The conjugate quaternion $q_m^{\ast}$ will achieve the opposite rotation.  Note that, in this section, we ignore the crystal symmetry of the individual grains; in other words, we consider the grains to have a structure described by point group \pg{1}.  The consequences of adding actual crystal symmetry, which leads to the concept of \indexit{disorientation}, are discussed elsewhere.

\begin{figure}[h!]
\centering\leavevmode
\includegraphics[width=0.75\textwidth]{\graphicspath GB.pdf}
\caption{\label{\chabbr:fig:GB}Simple symmetric tilt grain boundary in a cubic material. The $(100)$ plane is the boundary plane, and grains $1$ and $2$ are rotated in opposite sense by an equal angle $\theta$.}
\end{figure}

As an example, let us consider a planar grain boundary, as illustrated in Fig.~\ref{\chabbr:fig:GB}, where both grains have a single crystal axis in common, say they share the $\mathbf{e}_z$ axis, but they are rotated in opposite directions by an angle $\theta$ with respect to the boundary plane (this type of grain boundary is known as a \indexit{symmetric tilt boundary}). In the quaternion representation, we have for grain $1$ (with a counterclockwise rotation):
\[
	q_1 = \left[\cos\frac{\theta}{2},\left[0,0,\sin\frac{\theta}{2}\right]\right]\ ;
\]
and for grain $2$ we have:
\[
	q_2 = \left[\cos\frac{\theta}{2},\left[0,0,-\sin\frac{\theta}{2}\right]\right]\ ;
\]
For the misorientation quaternion we find:
\begin{equation*}
\begin{split}
	q_m &= \left[\cos^2\frac{\theta}{2}-\sin^2\frac{\theta}{2},\left[0,0,-2\cos\frac{\theta}{2}\sin\frac{\theta}{2}\right]\right]\ ;\\
	&=\left[\cos\theta,\left[0,0,-\sin\theta\right]\right]\ .
\end{split}
\end{equation*}
The misorientation angle $\Delta\theta$ is therefore an angle of magnitude $2\theta$.  Note that the rotation axis is $-\mathbf{e}_z$ since we have a clockwise rotation from grain $1$ to grain $2$.

One can show that the full quaternion multiplication is not necessary to obtain the misorientation angle (it is necessary, however, if the rotation axis is also needed). If we consider the quaternions as $4$-component vectors, then we can compute the dot product between those vectors using the normal approach; the dot product will then be equal to the cosine of half of the misorientation angle:
\[
	\cos\frac{\Delta\theta}{2} = q_1\cdot q_2 = \cos^2\frac{\theta}{2}-\sin^2\frac{\theta}{2} = \cos\theta\ ;
\]
note that we do \textit{not} use the conjugate quaternion $q_1^{\ast}$ in the dot product.  As before, we find for the misorientation angle that $\Delta\theta=2\theta$. 


\subsection{Uniform Orientation Grids}\label{\chabbr:sec:grids}
The cubochoric representation is useful because it allows for a two-step equal-volume mapping, via the homochoric representation, onto the Northern hemisphere, $\mathbb{S}_+^3$ of the unit quaternion sphere, which itself is a double-cover of $SO(3)$, the space of special orthogonal $3\times 3$ matrices.  If we create a uniform cubical grid of points in the cubochoric cell, each grid point will lie at the center of a small cube, and the volume of this cube will remain unchanged when mapped onto $\mathbb{S}_+^3$. The constant step size in the cubochoric grid will translate into a distribution of misorientation angles between neighboring grid points in $\mathbb{S}_+^3$; the main  reason for this distribution is the fact that an equal-volume mapping onto a (hyper)-sphere cannot also be an equal-angle mapping, so the constant grid step size in cubochoric space is mapped to a range of misorientation angles on the quaternion unit sphere.

As an example, consider a line of discrete equidistant orientations along the $c_z$ axis of the cubochoric cube with semi-edge length $a_c=\pi^{2/3}/2$; the coordinates can be  expressed as:
\[
	\mathbf{c}_j = (0,0, j \frac{a_c}{n})\text{ with }j\in[-n,n]\ .
\]
For this particular case we can apply the algorithm from  the previous section directly and we obtain for the homochoric coordinate:
\[
	\mathbf{h}_j = (\frac{3\pi}{4})^{1/3}\,(0,0, \frac{j}{n}) = \text{ with }j\in[-n,n] .
\]
We recognize the prefactor to be the radius of the homochoric ball $\mathbb{B}^3$; this means that, when $j=\pm n$, the homochoric point lies on the surface of the ball, which represents a rotation by an angle $\pi$.  All other rotations lie along the $h_z$ axis on a regularly spaced grid.  The conversion to quaternions requires numerical determination of the rotation angle $\omega$ for each homochoric grid point. Computation of the dot product between consecutive quaternions then allows for extraction of the misorientation angle, $\Delta\theta$.  Table~\ref{\chabbr:tb:cubogrid} lists all relevant quantities for a grid of $21$ (i.e., $n=10$) equidistant cubochoric points.

\begin{table}[h!]
\centering\leavevmode
{\small\begin{tabular}{cccccc}
\hline
$j$ & $c_z$ & $h_z$ & $q_0$ & $q_j\cdot q_{j+1}$ & $\Delta\theta_{j,j+1}$ \\
\hline\hline
-10 & -1.0725 & -1.3307 &  0.0000 &  0.9767 & 24.7723 \\
 -9 & -0.9653 & -1.1976 &  0.2145 &  0.9823 & 21.5969 \\
 -8 & -0.8580 & -1.0645 &  0.3937 &  0.9854 & 19.5723 \\
 -7 & -0.7508 & -0.9315 &  0.5442 &  0.9874 & 18.1836 \\
 -6 & -0.6435 & -0.7984 &  0.6699 &  0.9888 & 17.1934 \\
 -5 & -0.5363 & -0.6653 &  0.7734 &  0.9897 & 16.4763 \\
 -4 & -0.4290 & -0.5323 &  0.8562 &  0.9903 & 15.9603 \\
 -3 & -0.3218 & -0.3992 &  0.9197 &  0.9907 & 15.6023 \\
 -2 & -0.2145 & -0.2661 &  0.9645 &  0.9910 & 15.3761 \\
 -1 & -0.1073 & -0.1331 &  0.9911 &  0.9911 & 15.2664 \\
  0 &  0.0000 &  0.0000 &  1.0000 &  0.9911 & 15.2664 \\
  1 &  0.1073 &  0.1331 &  0.9911 &  0.9910 & 15.3761 \\
  2 &  0.2145 &  0.2661 &  0.9645 &  0.9907 & 15.6023 \\
  3 &  0.3218 &  0.3992 &  0.9197 &  0.9903 & 15.9603 \\
  4 &  0.4290 &  0.5323 &  0.8562 &  0.9897 & 16.4763 \\
  5 &  0.5363 &  0.6653 &  0.7734 &  0.9888 & 17.1934 \\
  6 &  0.6435 &  0.7984 &  0.6699 &  0.9874 & 18.1836 \\
  7 &  0.7508 &  0.9315 &  0.5442 &  0.9854 & 19.5723 \\
  8 &  0.8580 &  1.0645 &  0.3937 &  0.9823 & 21.5969 \\
  9 &  0.9653 &  1.1976 &  0.2145 &  0.9767 & 24.7723 \\
 10 & 1.0725 &  1.3307 &  0.0000 & --- & ---  \\
 \hline
\end{tabular}}
\caption{\label{\chabbr:tb:cubogrid}$z$ components of the cubochoric $(c_z)$ and homochoric $(h_z)$ points, the corresponding scalar part $q_0$ of the quaternion, the dot product between neighboring quaternions,  and the misorientation angle $\Delta\theta$ (in degrees) between points $j$ and $j+1$.}
\end{table}

The misorientation is nearly constant in the middle of the table and becomes larger at the start and end of the list. This table illustrates that an equal-volume sampling in the cubochoric cube produces an equal-volume sampling in both homochoric and quaternion representations, but the misorientation angle between consecutive sampling points is not constant. Note that the sum of the angles in the last column of the table is precisely $360^{\circ}$, as it should be since the line of grid points spans rotations from $-180^{\circ}$ to $+180^{\circ}$ around the $z$ axis.

Fig.~\ref{\chabbr:fig:planarsample}(a) shows the cubochoric cell with a uniform square grid of $9\times 9$ points in the $c_y$--$c_z$ plane; when those points are transformed to the stereographic representation (Fig.~\ref{\chabbr:fig:planarsample}(b)), they form circular sets of points; in other words, a grid of concentric squares in cubochoric space becomes a set of concentric circles in the stereographic representation. Since all neo-Eulerian representations share the property that the distance to the origin is equal to $f(\omega)$, this means that all points on a centered square in cubochoric space have the same rotation angle $\omega$. This remains true for a 3D uniform grid of concentric cubes which is mapped onto a 3D grid of concentric spheres.  The outer surface of the cubochoric cube represents all possible rotations by an angle $\pi$; note that only three facets of the cube (meeting in a single vertex) need to be taken into account to avoid double counting of $180^{\circ}$ rotations.

\begin{figure}[h!]
\centering\leavevmode
\includegraphics[width=0.9\textwidth]{\graphicspath planarsample.pdf}
\caption{\label{\chabbr:fig:planarsample}(a) uniform $9\times 9$ grid of points in the cubochoric $c_y$--$c_z$ plane; (b) representation of those points in a stereographic projection.}
\end{figure}

%%%%%%%%%%
% the following paragraph and figure have been moved to the OrientationSymmetry.tex chapter to avoid having a Foundational chapter reference a Topical chapter. [MDG, 12/31/25]

%When the cubochoric representation is combined with the Rodrigues representation, one can generate uniform orientation samples that are restricted to a crystallographic Fundamental Zone.  Consider the following two examples: the point group symmetries are \pg{24} and \pg{30}, with $N=20$ in the first case and $N=50$ in the second. Without the application of symmetry the total number of points generated would be $(2N+1)^3$ or $68,921$ for $N=20$ and $1,030,301$ for $N=50$.  The application of symmetry is performed by a simple test: first convert each cubochoric point to a Rodrigues vector, then test whether or not the point lies inside the RFZ; keep the point if it lies inside the RFZ, discard it otherwise.  The resulting lists of orientations have $n=5,367$ for \pg{24} and $n=41,625$ for \pg{30}.  The distributions of points along with the RFZ edges are shown as stereographic projections in Fig.~\ref{\chabbr:fig:samples}. Note in particular for \pg{30} that the sampling points lie on concentric spherical surfaces inside the RFZ due to the way they are sampled on concentric cubes in the cubochoric cell.
%
%\begin{figure}[h!]
%\centering\leavevmode
%\includegraphics[width=0.9\textwidth]{\graphicspath samples.pdf}
%\caption{\label{\chabbr:fig:samples}Stereographic projection of orientations uniformly sampled in cubochoric space and restricted to the Rodrigues Fundamental Zone for \pg{24} and \pg{30}.}
%\end{figure}
%%%%%%%%%%

\subsection{Random and Uniform Orientation Sampling}\label{\chabbr:ssec:randomsampling}
There are several algorithms available to generate random sets of orientations; in this section, we take a brief look at the most important approaches and make recommendations on whether or not they should be used. For each technique, we provide a figure that clearly illustrates any non-uniformities in the data set; the theory behind these maps, known as \indexit{square torus maps}, can be found in section~\ref{3DROTS:sec:CliffordTorus}. An additional method to evaluate uniformity of sampling on a hyper-sphere is explained in section~\ref{\chabbr:sec:Riesz}.  Each of the examples below uses a set of $5,000,000$ orientations generated by each of the sampling approaches.

\subsubsection{Naive uniform sampling}\label{\chabbr:sssec:naiveuniformSampling}
The simplest possible approach to generating a random set of orientations consists of randomly selecting four numbers uniformly on the interval $[-1,1]$. Most computer languages have a basic random number generator, usually on the interval $[0,1]$; such system-based generators generally have limited statistical quality. A better alternative for uniform sampling is the \indexit{Mersenne Twister} algorithm \cite{matsumoto1998a}. The four numbers are then collected as a quaternion, and the quaternion is normalized.  Fig.~\ref{\chabbr:fig:naive}{a} shows the square torus map for this sampling approach; the map clearly has non-uniformities due to the fact that the random uniform sample of four numbers does not intrinsically use any of the geometry or metric of the quaternion hyper-sphere.  

\begin{figure}[h!]
\centering\leavevmode
\includegraphics[width=0.66667\textwidth]{\graphicspath naive.pdf}
\caption{\label{\chabbr:fig:naive}$(X,Z_Y)$ square torus maps for (a) naive uniform sampling and (b) Marsaglia sampling.  Each map represents a sample of five million random orientations.}
\end{figure}

\subsubsection{Marsaglia sampling}\label{\chabbr:sssec:MarsagliaSampling}
In this approach \cite{marsaglia1972a} we uniformly draw two random numbers $a$ and $b$ from the interval $[-1,1]$ (using the Mersenne Twister); we repeat this step until the parameter $s_1=a^2+b^2\le 1$. Then we repeat this for an additional two numbers $c$ and $d$, until $s_2=c^2+d^2\le 1$. Then we replace $s_2$ by $\sqrt{(1-s_1)/s_2}$ and form the unit quaternion $q=[a,b,c\,s_2,d\,s_2]$. The resulting orientation set is more uniform than the set from the previous approach, but there is still some non-uniformity that shows up as vertical intensity bands in Fig.~\ref{\chabbr:fig:naive}{b}; the level of uniformity within the central vertical band is reasonable.

\subsubsection{Shoemake sampling}\label{\chabbr:sssec:ShoemakeSampling}
The Shoemake algorithm \cite{shoemake1992a} relies on three uniformly selected random numbers, $u_1$ on the interval $[0,1]$, and $u_2$ and $u_3$ from the interval $[0,2\pi]$. The Mersenne Twister algorithm is most suitable to generate these numbers.  Setting $a=\sqrt{1-u_1}$ and $b=\sqrt{u_1}$, the random unit quaternion is then obtained as $q=[a\sin u_2, a\cos u_2, b \sin u_3, b\cos u_3]$.  This approach generates a highly uniform set of orientations, as evidenced qualitatively by the square torus map in Fig.~\ref{\chabbr:fig:shoemake}(a).

\subsubsection{Cubochoric sampling}\label{\chabbr:sssec:CubochoricSampling}
This approach does not rely on random numbers, but instead takes a uniform cubical grid in the cubochoric cube and transforms all the points to unit quaternions via the homochoric representation \cite{degraef2014r}. Since the unit quaternions only have a positive scalar part in this approach, the corresponding square torus map only shows the central vertical band (Fig.~\ref{\chabbr:fig:shoemake}(b)). One difference between this technique and the next one compared to the previous three approaches is that the cubochoric sampling generates a uniform orientation set on $\mathbb{S}^3_+$ instead of a random set.

\begin{figure}[h!]
\centering\leavevmode
\includegraphics[width=\textwidth]{\graphicspath shoemake.pdf}
\caption{\label{\chabbr:fig:shoemake}$(X,Z_Y)$ square torus maps for (a) Shoemake sampling, (b) cubochoric sampling, and (c) super-Fibonacci sampling.  Each map represents a sample of five million random orientations.}
\end{figure}

\subsubsection{Super-Fibonacci sampling}\label{\chabbr:sssec:FibonacciSampling}
Super-Fibonacci sampling was introduced in 2022 \cite{alexa2022a} as a fast sampling technique based on two irrational numbers, $(\phi,\psi)$.  There are several choices for these numbers and here we set $\phi$ equal to the \indexit{golden ratio}, $\phi=(1+\sqrt{5})/2=1.618033989\ldots$, and $\psi$ is the \indexit{super-golden ratio}, the positive root of the equation $\psi^3=\psi^2+1$ so that:
\[
	\psi = \frac{1}{3}+\frac{2}{3}\cos\left(\frac{1}{3}\arccos\left(\frac{29}{2}\right)\right)=1.4655712318\ldots\ .
\]
To generate a uniformly distributed set of $N$ unit quaternions, we employ the index $i\in[0\ldots N-1]$ and compute the following series of parameters:
\begin{equation*}
\begin{split}
	s &= i+\frac{1}{2}\rightarrow d = 2\pi\,s\ ;\\
	t &= \frac{s}{N}\rightarrow r = \sqrt{t}\text{ and }R = \sqrt{1-t}\ ;\\
	\alpha &= \frac{d}{\phi}\text{ and }\beta = \frac{d}{\psi}\ .
\end{split}
\end{equation*}
The unit quaternions are then given by $q_i=[r\sin\alpha, r\cos\alpha, R\sin\beta, R\cos\beta]$. The square torus map for this sampling is shown in Fig.~\ref{\chabbr:fig:shoemake}(c), and is remarkably featureless.  This approach is simple, fast, and easy to implement and produces a highly uniform set of orientations.  

%\subsubsection{Other sampling approaches}\label{\chabbr:sssec:OtherSampling}


\section{Evaluation of Sampling Uniformity by Means of Riesz Kernels}\label{\chabbr:sec:Riesz}
In this section\footnote{This section contains more advanced math and may be skipped on first reading.}, we take a brief look at what sampling uniformity really means.  One of the main approaches to the evaluation of uniformity dates back to J.J. Thomson, who asked the question: if we have a set of $N$ electrons and we place them on the surface of a conducting sphere, then what is the configuration that minimizes the electrostatic energy? This is known as the \indexit{Thomson problem}, and it has analytical solutions for only a few small values of $N$, for which the configurations are mostly the vertices of regular polyhedra.  This problem can be generalized to higher dimensional spheres, including the quaternion unit sphere, and from a mathematical point of view, the electrostatic interaction energy is replaced by the \indexit{Riesz kernel}, which is defined in this section.  

The Thomson problem is based on the standard electrostatic interaction energy between charges, which is given by:
\[
	U_{ij} = \frac{1}{4\pi\epsilon_0}\frac{q_iq_j}{r_{ij}}\ ,
\]
where $\epsilon_0$ is the permittivity of vacuum and $r_{ij}$ is the Euclidean distance $\vert\mathbf{r}_i-\mathbf{r}_j\vert$ between the two charges. Multiplying both sides by $4\pi\epsilon_0$ and summing over all charge pairs (assuming electrons) we have for the total reduced Coulomb interaction energy for $N$ charges, subject to the constraint that $\vert\mathbf{r}_i\vert=1$ for all $i$:
\[
	\bar{U}(N) = \sum_{i<j} \frac{1}{r_{ij}}\ .
\]
There is no general analytical solution for this problem, so minimization is typically performed by numerical algorithms. The function $1/r$ can be replaced by any other monotonically increasing function; a common choice is $f(r)=r^{-s}$, with $s$ an integer between $1$ and the dimensionality $d$ of the hyper-sphere, which is known as the \indexit{Riesz $s$-kernel}.  The generalized expression for the total Coulomb interaction energy becomes:
\[
	\bar{U}(N;s) = \sum_{i<j} \frac{1}{r^s_{ij}}\ .
\]

There is a tremendous amount of mathematical literature about this kind of generalized energy sums, in particular for higher-dimensional spaces.  This problem is relevant for orientation sampling, since unit quaternions live on the three-sphere $\mathbb{S}^3$; the lower the \indexit{Riesz energy} for a given set of orientations, the more uniform the orientation sample is.  In  the case of unit quaternions, we have $d=3$, so that the possible values for $s$ are $1$, $2$, and $3$.  We will represent the Riesz energy for a set $\mathcal{S}_M$ of $M$ unit quaternions as:
\begin{equation}
	E_s(\mathcal{S}_M)\equiv \sum_{i\ne j} \vert\vert q_i-q_j\vert\vert^{-s}=2\sum_{1\le i< j\le M} \vert\vert q_i-q_j\vert\vert^{-s}\ ; \label{\chabbr:eq:Riesz}
\end{equation}
the symbol $\vert\vert\cdot\vert\vert$ represents the Euclidean distance in $\mathbb{R}^4$ (since we are working on a hyper-sphere, this is also known as the \indexit{chordal distance}). This distance is computed in the standard way by considering the pairwise differences between the quaternion components $q_{i,k}$ with $k\in[0\ldots 3]$:
\[
	\vert\vert q_i-q_j\vert\vert = \sqrt{\sum_{k=0}^{3}(q_{i,k}-q_{j,k})^2}.
\]
It should be noted that the summation range $1\ldots M$ includes both antipodal quaternions $q_i$ and $-q_i$, so that $M=2N$ where $N$ is the number of orientation quaternions with positive scalar part $q_{i,0}$.  

The optimal Riesz energies for perfectly uniform covering of the three-sphere $\mathbb{S}^3$ with $M$ quaternions (including antipodal points) are known and given by:
\begin{equation}
\begin{split}
	E_1^{\text{opt}}(\mathcal{S}_M) &= \frac{8}{3\pi}M^2\ ; \\
	E_2^{\text{opt}}(\mathcal{S}_M) &= M^2\ ; \\
	E_3^{\text{opt}}(\mathcal{S}_M) &= \frac{2\log M}{3\pi}M^2\ .
\end{split}
\end{equation}
It has been established \cite{hardin2005a} that these are the leading terms in series expansions in terms of the parameter $M^{1+s/3}$;  it is known that the next term in the expansion is positive.  Therefore, computing the energies $E_s(\mathcal{S}_M)$ using eq.~(\ref{\chabbr:eq:Riesz}) and comparing them with the analytical results as a ratio:
\[
	r_s \equiv \frac{E_s(\mathcal{S}_M)}{E_s^{\text{opt}}(\mathcal{S}_M)}\ ,
\]
results in a simple test for sampling uniformity: \textit{the lower the ratios $r_s$ for a given set $\mathcal{S}_M$, the more uniform the sampling.} It is possible for the ratios to drop slightly below unity due to the unknown higher order terms in the expansions.

As an example of this test, we list in Table~\ref{\chabbr:tb:Riesz} the ratios $r_s$ for a sampling of $M=200,000$ orientations, using the various sampling techniques described in section~\ref{\chabbr:ssec:randomsampling}.  Note that most values are close to or just below $1$; the value of $r_3$ is sensitive to local non-uniformities and clearly ranks the uniformity of the sampling techniques.  The number of orientations in the Riesz computation for cubochoric sampling is of the form $M=2(2N+1)^3$, in this case with $N=23$ to arrive at a number close to $M=200,000$, which was used for the other sampling techniques.

\begin{table}[h]
\centering\leavevmode
\begin{tabular}{|l|r|r|r|}
\hline
Sampling Method &   $r_1$   &   $r_2$ & $r_3$ \\
\hline
Uniform &                  1.00724507 &    1.05284771 &     2.68756411 \\
Marsaglia &               1.01138374 &    1.05208453 &     2.19656912 \\
Shoemake &              0.99999244 &    1.00008561 &    2.08396706 \\
Cubochoric &             0.99974563 &    0.98480548 &    1.02266132 \\
Super-Fibonacci &     0.99968683 &    0.98198723 &    0.97808663 \\
\hline
\end{tabular}
\caption{\label{\chabbr:tb:Riesz}Normalized Riesz energies $r_i$ for five different sampling approaches; each sample contained $M=200,000$ orientations, except for the Cubochoric sample which contained $194,672$ orientations.}
\end{table}

%\subsection{Dispersion}\label{\chabbr:ssec:dispersion}
%
%

\section{Orientations and Directional Statistics}\label{\chabbr:sec:orientationaveraging}
The topic of averaging data on a circle or on a sphere belongs to the branch of statistics that is known as \indexit{Directional Statistics} (DS); DS defines how data on a circle (i.e., data that is periodic) should be treated when averages or other statistical parameters are needed. Similarly, data on spheres and hyper-spheres requires special handling in the context of averaging.  Since 3D rotations and orientations can be represented on the hyper-sphere $\mathbb{S}^3$, it stands to reason that the techniques and concepts introduced in this section will be important for the study of material textures.  We begin with data that lives on a circle.

\subsection{Averaging of Circular Data}\label{\chabbr:ssec:circleaveraging}

\begin{floatingfigure}[r]{3.0in}
\centering
{\includegraphics[width=2.8in]{\graphicspath winddirections.pdf}}
{\caption{\label{\chabbr:fig:winddirections}\small Made-up data set representing the daily average wind directions for Pittsburgh, PA, in $5^{\circ}$ bins from North.}}
\end{floatingfigure}
\noindent By its very nature, a circle is a periodic shape, and functions defined on the circle are periodic. As an example, consider the function $f(\theta) = \mathrm{e}^{\mathrm{i}\theta}$ which can be written as $\cos\theta+\mathrm{i}\sin\theta$ by Euler's Formula.  $f(\theta)$ is a unit modulus complex number and it is clear from the properties of the trigonometric functions that $f(\theta)=f(\theta+2\pi\,k)$ with $k\in\mathbb{Z}$.  As an example (using made-up data) we consider the daily average wind direction for Pittsburgh, PA, binned in $5^{\circ}$ bins and measured in degrees from North, as shown in Fig.~\ref{\chabbr:fig:winddirections}. 

Since this is periodic data, we can represent it as points on a circle, as shown in Fig.~\ref{\chabbr:fig:windcircle}(a) or by converting the linear histogram into a polar histogram as shown in Fig.~\ref{\chabbr:fig:windcircle}(b). If we naively compute the average direction as a standard numerical average over the range $[0^{\circ}\ldots 360^{\circ}]$ then we obtain $\langle\theta\rangle=107.06^{\circ}$, which is indicated by the red radial line in  Fig.~\ref{\chabbr:fig:windcircle}(b).  

\begin{figure}[h!]
\centering\leavevmode
\includegraphics[width=0.8\textwidth]{\graphicspath windcircle.pdf}
\caption{\label{\chabbr:fig:windcircle}Representation of the wind direction data on a circular plot (a) and a polar plot (b). Colored lines represent the average directions for two different angular ranges.}
\end{figure}

If we change the horizontal axis of the histogram in Fig.~\ref{\chabbr:fig:winddirections} to the range $[-180^{\circ}\ldots 180^{\circ}]$, and wrap the data periodically, then we obtain for the average angle $\langle\theta\rangle=52.81^{\circ}$, represented by the blue line. 
So, the value of the conventional average depends on where we pick the starting point on the circle!  From this we learn that \textit{we should never simply take the average of a series of angles or directions}; in particular, one should \textit{never} compute the average of a set of Euler angle triplets by simply taking the average for each angle separately.

For the correct procedure to compute the average for circular data we turn to Directional Statistics. We begin by treating the data as a collection of unit vectors in a 2D plane; the vectors are defined as
\[
	\hat{\mathbf{r}}_i = (\cos\theta_i,\sin\theta_i)\ .
\]
The center-of-mass of these vectors is given by the vector $(\bar{C},\bar{S})$ with 
\begin{equation*}
\begin{split}
	\bar{C} &= \frac{1}{n}\sum_{i=1}^n \cos\theta_i\ ;\\
	\bar{S} &= \frac{1}{n}\sum_{i=1}^n \sin\theta_i\ .	
\end{split}
\end{equation*}
Note that we do use the traditional expression for the average, but we apply it to the vector components, not to the angles.  The center-of-mass vector can be written in polar form as: $(\bar{C},\bar{S})=\bar{R}(\cos\bar{\theta},\sin\bar{\theta})$ where $\bar{R}\equiv \sqrt{\bar{C}^2+\bar{S}^2}$ is the \indexit{mean resultant length}, and $\bar{\theta}$ is the mean direction, defined as:
\begin{equation*}
	\bar{\theta} = \left\{\begin{array}{lcl}
	\text{arctan}(\bar{S}/\bar{C}) & \qquad & \bar{C}\ge 0\\
	\text{arctan}(\bar{S}/\bar{C})+\pi & \qquad & \bar{C}< 0\end{array}\right.\ .	
\end{equation*}
When we apply these relations to the wind directions data set, we obtain $\bar{C}=0.539897$ and $\bar{S}=0.345104$, which translates to $\bar{R}=0.640750$ and $\bar{\theta}=32.588^{\circ}$.  Since the angle for the vectors is measured starting from the $x$-axis, and the wind directions are with respect to North, we take the complement of this angle and find $\bar{\theta}=57.41^{\circ}$, which is indicated by the green line in  Fig.~\ref{\chabbr:fig:windcircle}(b).  

The mean resultant length $\bar{R}$ is a measure for how ``concentrated'' the data is; the larger $\bar{R}$, the narrower the distribution of values is. When we introduce distributions on spheres and hyper-spheres in the next section, the parameter corresponding to $\bar{R}$ will be called the \indexit{concentration parameter}.  Other statistical parameters for the circular data case include:
\begin{itemize}
	\item the \indexit{sample circular variance}: $V=1-\bar{R}$ ($=0.359$ for the wind direction data)\ ;
	\item the \indexit{circular standard deviation}: $v=(-2\ln\bar{R})^{1/2}$ ($=0.943$)\ ;
	\item the \indexit{median direction}: represent the data on the range $[0^{\circ}\ldots 360^{\circ}]$; shift the data to the left by an angle $\theta_m$ such that half of the data points have negative angles and half positive (for the wind directions data set we have $\theta_m=75.40^{\circ}$)\ ;
	\item for the definition of other circular statistics quantities, including circular dispersion, circular mean difference, and circular range, we refer the interested reader to section 2.3 in \cite{mardia2009a}. 
\end{itemize}

The example above dealt with discrete data; it is also possible to work with continuous distribution functions. We define a distribution function $F(x)$ as the probability $\text{Pr}(0<\theta\le x)$ that a randomly selected angle $\theta$ lies in the stated range, with $0\le x\le 2\pi$.  For circular data this function has the property that $F(x+2\pi)-F(x)=1$ for all $x$ values.  The probability that the angle $\theta$ falls in the range $[\alpha,\beta]$ is then given by:
\[
	\text{Pr}(\alpha<\theta\le\beta)=F(\beta)-F(\alpha)=\int_{\alpha}^{\beta}f(\theta)\,\mathrm{d}\theta\ ,
\]
where $f(\theta)$ is the probability density function; this function is positive-valued, periodic, and normalized.  As a simple example, for a uniform probability density, $f(\theta) = 1/2\pi$, we have:
\[
	\text{Pr}(\alpha<\theta\le\beta)= \frac{\beta-\alpha}{2\pi}\ .
\]
A more interesting and useful distribution is the so-called \indexit{von Mises Distribution}, which can be regarded as the circular equivalent of the normal distribution on the 1D line.  The von Mises distribution $M(\mu,\kappa)$, where $\mu$ is the mean direction and $\kappa$ the concentration parameter, has probability density function:
\begin{equation}
	g(\theta;\mu,\kappa)\equiv \frac{1}{2\pi\,I_0(\kappa)}\mathrm{e}^{\kappa\cos(\theta-\mu)}\ ;
\end{equation}
In this expression, $I_0(\kappa)$ is a modified Bessel function of the first kind and order $0$, which can be described by the series expansion:
\[
	I_0(x) = \sum_{k=0}^{\infty}\frac{1}{(k!)^2}\left(\frac{x}{2}\right)^{2k}=1+\frac{1}{4}x^2+\frac{1}{64}x^4+\frac{1}{2,304}x^6 + \frac{1}{147,456}x^8 + \ldots\ .
\]

\begin{figure}[h!]
\centering\leavevmode
\includegraphics[width=0.9\textwidth]{\graphicspath vonMises1D.pdf}
\caption{\label{\chabbr:fig:vonMises1D}(a) plots of the von Mises probability density function $g(\theta;\mu,\kappa)$ for $\mu=0$ and different values of the concentration parameter $\kappa$; (b) polar plot for $\mu=\pi/6$ and $\kappa=8$.}
\end{figure}

Fig.~\ref{\chabbr:fig:vonMises1D}(a) shows the von Mises probability density function $g(\theta;0,\kappa)$ for $\theta\in[-180^{\circ},180^{\circ}]$ and four different values of $\kappa$; note that this function becomes more narrowly peaked as the concentration parameter increases. Fig.~\ref{\chabbr:fig:vonMises1D}(b) shows the function $g(\theta;\pi/6,8)$ as a polar plot.  


\subsection{Distributions on (hyper)-spheres}\label{\chabbr:ssec:sphereaveraging}

\subsubsection{Discrete data sets}
While the theory described in this section is valid for data in higher dimensional spaces of arbitrary dimension $p$, we will focus here on the $3$-sphere $\mathbb{S}^3$ in $\mathbb{R}^4$, so that $p=4$.  By analogy with the circle case, we will not consider the angles directly, but instead we represent data points by unit vectors $\mathbf{x}_i$ pointing to points on $\mathbb{S}^3$; think of this data set as a collection of grain orientations represented as unit quaternions (but pay attention to the ``word of caution'' at the end of this section).  As in the circle case, we take the conventional average over the entire sample of data points, which leads to the sample mean vector 
$\bar{\mathbf{x}}$:
\begin{equation}
	\bar{\mathbf{x}} = \frac{1}{n}\sum_{i=1}^n \mathbf{x}_i\ .\label{\chabbr:eq:regularaverage}
\end{equation}
While the $\mathbf{x}_i$ are unit vectors (unit quaternions), the mean vector is not a unit vector so it falls in the space $\mathbb{R}^4$.  We can, however, express the mean vector in terms of its length, the \indexit{mean resultant length} $\bar{R}$, and the \indexit{mean direction vector} $\bar{\mathbf{x}}_0$: 
\[
	\bar{\mathbf{x}} = \bar{R}\, \bar{\mathbf{x}}_0\ .
\] 
For any data set of vectors $\mathbf{x}_i$ we can ask the question: is this data strongly clustered or is it uniform?  One approach to answering this question is by means of the \indexit{scatter matrix} $\mathbf{T}$, which is defined by means of the outer products $\mathbf{x}_i\otimes \mathbf{x}_i$:\footnote{The outer product of two vector $\mathbf{u}$ and $\mathbf{v}$ is defined in terms of the components by $T_{ij} = u_iv_j$.}
\[
	\mathbf{T} = \sum_{i=1}^n \mathbf{x}_i\otimes \mathbf{x}_i\ .
\]
For $p=4$, this is a $4\times 4$ matrix; the nature of a distribution on a sphere can be determined from the four eigenvalues $t_i$ of this matrix. If we rank the eigenvalues from smallest, $t_1$, to largest, $t_4$, then the following holds: 
\begin{itemize}
\item if all four eigenvalues are approximately equal to each other, then the data is distributed uniformly on the sphere;
\item if $t_4$ is significantly larger than the other three, and $t_1$, $t_2$, and $t_3$ are approximately equal, then the distribution is clustered and has rotational symmetry around eigenvector $\mathbf{t}_4$; if in addition the mean resultant length $\bar{R}$ is close to unity, then the distribution is unimodal (single peak).
\item if $t_1$ is significantly smaller than the other three, and those are approximately equal, then we have a symmetric ``girdle'' distribution in the equatorial plane normal to eigenvector $\mathbf{t}_1$.
\end{itemize}
For more information about discrete data sets on hyper-spheres we refer the interested reader to section 9.2 in \cite{mardia2009a}.  

One \textbf{word of caution} is needed at this point: the averaging procedure described above is valid for a general three-sphere $\mathbb{S}^3$ in $\mathbb{R}^4$. From a texture point of view, we use the three-sphere to represent 3D orientations using unit-quaternions, and in that case there is an equivalence between $q$ and $-q$, i.e., $\mathbb{S}^3$ is a double-cover of the rotation group. If we take a set of $n$ rotation quaternions $q_i$, and we replace half of them by $-q_i$, then this will clearly change the average defined in eq.~\ref{\chabbr:eq:regularaverage}.  In particular, if we included $q_i$ and $-q_i$ for all orientations, the average would always be zero.  It is therefore important to work only with orientations that lie in a hemisphere of $\mathbb{S}^3$, and it is customary to require that the scalar part of the quaternion be positive or zero; we can conveniently call this hemisphere the ``Northern hemisphere'' of $\mathbb{S}^3$ and denote it by $\mathbb{S}^3_+$.  In section~\ref{\chabbr:sec:oraverages3} we will describe other procedures for rotation quaternion averaging.

\subsubsection{Continuous distributions on hyper-spheres}%verify the normalization parameters !!!
The von Mises distribution on the circle can be generalized to (hyper-)spheres in which case we talk about \indexit{spherical data distributions}. In  $p$-dimensional space, the \indexit{von Mises-Fisher} (vMF) probability density is defined by:
\begin{equation}
	f(\mathbf{x};\bm{\mu},\kappa)\equiv \frac{\kappa^{\frac{p}{2}-1}}{(2\pi)^{p/2}I_{\frac{p}{2}-1}(\kappa)}\mathrm{e}^{\kappa\bm{\mu}\cdot\mathbf{x}}\ ,
\end{equation}
where, as before, $\bm{\mu}$ is the mean direction and $\kappa$ the concentration parameter; $\mathbf{x}$ is a random vector in the $p$-dimensional space $\mathbb{R}^p$. In the denominator we have the Euler gamma function $\Gamma(x)$ and the modified Bessel function of the first kind and order $\frac{p}{2}-1$.  For quaternions, we have $p=4$, so the vMF probability density simplifies to:
\begin{equation}
	f(q;\mu,\kappa) \equiv \frac{\kappa}{4\pi^2 I_1(\kappa)}\mathrm{e}^{\kappa\mu\cdot q}\ ;\label{\chabbr:eq:vMFpdf}
\end{equation}
in this expression, the exponential contains the standard dot product between two quaternions, the mean direction $\mu$ and the random orientation $q$, as four-component vectors in $\mathbb{R}^4$.  One can show that this distribution is rotationally symmetric around $\mu$; if $\kappa=0$ then the distribution is uniform, if $\kappa>0$ then the orientations are clustered around $\mu$.  The expectation value $E[q]$ of this distribution is $\rho\mu$ with $\rho=I_2(\kappa)/I_1(\kappa)$, i.e., the expectation value is proportional to the mean direction.  Fig.~\ref{\chabbr:fig:vMFW}(a) shows the vMF probability density for several values of $\kappa$; the dot product $\mu\cdot q$ was replaced by $\cos\theta$ with $\theta\in[0,\pi]$.  As for the circle case, the vMF probability density can be considered as the spherical equivalent of the normal distribution.

\begin{figure}[h!]
\centering\leavevmode
\includegraphics[width=0.9\textwidth]{\graphicspath vMFW.pdf}
\caption{\label{\chabbr:fig:vMFW}(a) plots of the von Mises-Fisher probability density function $f(q;\mu,\kappa)$ for $\mu\cdot q=\cos\theta$ and different values of the concentration parameter $\kappa$; (b) plots of the Watson distribution for the same parameters.}
\end{figure}

When considering orientation distributions in highly textured materials, the concentration parameter $\kappa$ can become large; values above $1,000$ are not uncommon.
Accurate computation of the prefactor becomes problematic when using the standard numerical implementation of the modified Bessel function $I_1(\kappa)$; it can be shown that, for $\kappa>30$, the following expression is a highly accurate approximation of the vMF density function:
\begin{equation}
	f(q;\mu,\kappa) \approx \frac{512}{\sqrt{2}\pi^{3/2}}\frac{\mathrm{e}^{\kappa(\mu\cdot q-1)}\kappa^{9/2}}{1024\kappa^3-384\kappa^2-120\kappa-105}\ .	
\end{equation}

The \indexit{Watson probability density} explicitly takes into account the equivalence between the quaternions $q$ and $-q$ by squaring the argument of the exponential. For $p=4$, the Watson function is defined by:
\begin{equation}
	f_W(q;\mu,\kappa) \equiv \frac{\mathrm{e}^{-\frac{\kappa}{2}}}{I_0(\frac{\kappa}{2})-I_1(\frac{\kappa}{2})}\mathrm{e}^{\kappa(\mu\cdot q)^2}\ ;\label{\chabbr:eq:Watsonpdf}
\end{equation}
This function is shown in Fig.~\ref{\chabbr:fig:vMFW}(b) and has two peaks, one at $\theta=0$, corresponding to $\mu$, and one at $\theta=\pi$, corresponding to $-\mu$. This function also has rotational symmetry around the $(\mu,-\mu)$ axis and therefore represents axial data. As for the vMF density, the prefactor of the Watson density function becomes numerically problematic when $\kappa>30$; an accurate approximation for large $\kappa$ values is given by:
\begin{equation}
	f_W(q;\mu,\kappa) \approx 128\sqrt{\pi}\frac{\mathrm{e}^{\kappa((\mu\cdot q)^2-1)}\kappa^{9/2}}{128\kappa^3+96\kappa^2+180\kappa+525}\ .	
\end{equation}
For a given value of $\kappa$, the Watson pdf tends to give rise to a tighter orientation distribution than the von Mises-Fisher pdf.

Both vMF and Watson probability density functions (pdfs) have rotational symmetry around the direction $\mu$. If we consider the $n$-dimensional normal pdf:
\begin{equation}
	f(x_1,x_2,\ldots,x_n)=\frac{1}{\sqrt{(2\pi)^n\vert\bm{\Sigma}\vert}}\mathrm{e}^{-\frac{1}{2}(\mathbf{x}-\bm{\mu})^T\bm{\Sigma}^{-1}(\mathbf{x}-\bm{\mu})}\ ,
\end{equation}
where $\bm{\mu}$ is the mean (or mode) and $\bm{\Sigma}$ is the \indexit{covariance matrix}, then this pdf will have rotational symmetry when the covariance matrix is a diagonal matrix with the same value for each entry. For the 2D case, this means that a projection of the function onto the $(x_1,x_2)$ coordinate plane will have circular symmetry. When the covariance matrix has different values along the diagonal, then the circular projection will become elliptical, with the main axes of the ellipse aligned with the coordinate axes. For a general covariance matrix, the elliptical projections will have principal axes that do not align with the coordinate axes.  

A similar situation arises for spherical distributions; the \indexit{Bingham probability density} function must be used when the data does not have rotational symmetry around the mean direction.  In this general case, there is no single value for $\kappa$, and $\kappa$ and $\mu$ can not be separated from each other in the mathematical formalism. The Bingham pdf for $p=4$ is defined as:
\begin{equation}
	f_B(\pm q;\mathbf{A})\equiv \frac{1}{{}_1F_1\left(\frac{1}{2},2,\mathbf{A}\right)}\mathrm{e}^{q^T\mathbf{A}q}\ ,
\end{equation}
where $\mathbf{A}$ is a traceless $4\times 4$ matrix and the function ${}_1F_1(a,b;z)$ is a hypergeometric function defined as:
\[
	{}_1F_1(a,b;z) = \sum_{k=0}^{\infty}\frac{(a)_k}{k!(b)_k}z^k\text{ with } (d)_k\equiv\frac{\Gamma(d+k)}{\Gamma(d)}\ .
\]
The symbol $(d)_k$ is known as a \indexit{Pochhammer symbol}.  Note that, for the Bingham pdf, the argument $z$ of the hypergeometric function is a $4\times 4$ matrix, so the computation of the normalization prefactor represents a non-trivial numerical problem.  

\subsection{Orientation sampling on $\mathbb{S}^3$}\label{\chabbr:ssec:orsamplings3}
Before we can describe sampling on $\mathbb{S}^3$, we need to introduce two important concepts: the \indexit{tangent-normal decomposition}, and \indexit{rejection sampling}.  Then we will describe random sampling from the von Mises-Fisher and Watson probability density functions.

\subsubsection{The tangent-normal decomposition}
In this section, we describe how the \indexit{tangent-normal decomposition} is defined for points on $\mathbb{S}^3$; the decomposition is much more general and can be defined on arbitrary surfaces in arbitrary dimensions, but for our purposes the spherical case in $\mathbb{R}^4$ will suffice. Consider a general point $\mathbf{x}$ on $\mathbb{S}^3$; we pick a \textit{reference direction} $\bm{\mu}$, also on the sphere, and determine the \indexit{tangent space} at the point $\bm{\mu}$ (see Fig.~\ref{\chabbr:fig:tndecomp}(a) for a 2D representation).  One can then decompose the vector $\mathbf{x}$ into components along $\bm{\mu}$ and along the unit tangent vector $\bm{\xi}$. We find:
\begin{equation}
	\mathbf{x}= t\,\bm{\mu} + \sqrt{1-t^2}\,\bm{\xi}= (\mathbf{x}\cdot\bm{\mu})\,\bm{\mu} + \sqrt{1+(\mathbf{x}\cdot\bm{\mu})^2}\,\bm{\xi}\ ,
\end{equation}
where in the second equality we have made use of the fact that $t=\mathbf{x}\cdot\bm{\mu}$. The reference direction $\bm{\mu}$ can be chosen arbitrarily, but the choice of symbol suggests that one often uses the mean direction of a distribution as the reference direction. 

\begin{figure}[h!]
\centering\leavevmode
\includegraphics[width=0.6\textwidth]{\graphicspath tndecomp.pdf}
\caption{\label{\chabbr:fig:tndecomp}(a) 2D illustration of the tangent-normal decomposition; (b) Illustration of the possible outcomes of the Buffon needle experiment (gray rectangle) along with the region below the sinusoid which represents the successful outcomes.}
\end{figure}


\subsubsection{Rejection sampling}
The concept of the rejection sampling algorithm is very similar to the way Buffon's Needle experiment is modeled. In this experiment, a needle of length $L$ is randomly thrown onto a sheet of paper onto which parallel lines with spacing $D$ have been drawn. A random sample (i.e., a single toss of the needle) is rejected if the needle does not cross any of the lines, and it is accepted if one (or more) lines are intercepted. If the orientation of the needle is described by its angle $\theta\in[0,\pi]$ with respect to the lines, then it is easy to see that the needle will intersect a line if the distance $y$ between one end of the needle and the nearest line is less than or equal to $L\sin\theta$. The maximum value of $y$ is $D$, so that all attempts will produce points in a rectangle of area $D\pi$; all successful attempts will fall on or below the curve $L\sin\theta$, so that the probability density function for this geometry is
\[
	f(\theta) = \frac{L\sin\theta}{\pi D}\ .
\]
The probability of intersecting a line is then:
\[
	P = \frac{L}{\pi D}\int_0^{\pi}\sin\theta\,\mathrm{d}\theta = \frac{2L}{\pi D}\ .
\]
Rejection sampling in this case can be described as follows (see Fig.~\ref{\chabbr:fig:tndecomp}(b)): consider a rectangular board with dimensions $\pi\times D$; on the board the curve $y=L\sin\theta$ is drawn. If a large number $N$ of darts is randomly thrown at the board, and a total of $C$ darts fall on or below the curve, then $P=C/N$. The attempts above the curve are rejected, those below the curve are accepted.  This approach can be used to randomly sample any distribution.

\subsubsection{Random sampling of a von Mises-Fisher distribution}
To generate a random sample of orientations for a vMF distribution with given mean direction $\bm{\mu}$ and concentration parameter $\kappa$, we combine the concepts of tangent-normal decomposition and rejection sampling:
\begin{enumerate}
	\item Use normal random sampling to select a random 3D unit vector $\mathbf{u}$ on $\mathbb{S}^2$; this will become the tangent space unit vector $\bm{\xi}$ of the tangent-normal decomposition by construction of the unit quaternion $[0,\mathbf{u}]$. We select $[1,0,0,0]$ as the initial normal vector, which in a later step we will rotate to the desired direction $\bm{\mu}$.
	\item Use rejection sampling to select a random number $t\in[0,1]$ from the distribution:
	\begin{equation}
		f(t) = \frac{\kappa^{\frac{5}{2}}\sqrt{1-t^2}\mathrm{e}^{\kappa t}}{\sqrt{2\pi}[\kappa\cosh(\kappa)-\sinh(\kappa)]}\ .\label{\chabbr:eq:rejectionnMF}
	\end{equation}
	This function is obtained by integrating the pdf around its mean direction $\bm{\mu}$.
	\item Combine $t$ and $\mathbf{u}$ using the tangent-normal decomposition into the quaternion:
	\begin{equation}
		q = [t,\mathbf{0}] + \sqrt{1-t^2}[0,\mathbf{u}]\ .
	\end{equation}
	By construction this is a unit quaternion; also, the tangent vector $[0,\mathbf{u}]$ is perpendicular to the normal vector $[1,\mathbf{0}]$, as is required for the tangent-normal decomposition.
	\item Rotate the resulting quaternion $q$ to the desired mean direction $\bm{\mu}$ (see description below).
	\item Repeat this process until the desired number of samples is reached.
\end{enumerate}

If sampling from the Watson pdf is needed, the same process can be followed but the rejection sampling function $f(t)$ must be replaced by:
\begin{equation}
	f(t) = \frac{2}{\pi}\frac{\mathrm{e}^{-\frac{\kappa}{2}}}{I_0(\frac{\kappa}{2})-I_1(\frac{\kappa}{2})}\sqrt{1-t^2}\mathrm{e}^{\kappa t^2}\ .
\end{equation}
Since both the vMF and Watson pdfs have rotational symmetry around their respective mean directions $\bm{\mu}$, one can relate the concentration parameter $\kappa$ to an angular interval $\Delta\theta$, which can be regarded as being equivalent to a full-width-at-half-maximum for a normal distribution.  The relation between $\kappa$ and $\Delta\theta$ for both pdfs is given by:
\begin{equation}
	\Delta\theta = \arccos(1-\frac{1}{\kappa})\ .
\end{equation}


The rotation matrix needed to rotate all generated quaternions to the mean direction is a $4\times 4$ rotation matrix, i.e., a rotation in $\mathbb{R}^4$. A $4\times 4$ matrix, $A$, is initialized with all zeroes, and the components of the mean direction $\bm{\mu}$ are copied in the first column of $A$. A numerical routine\footnote{For instance, in Matlab this would be the null() operation; in the LAPACK numerical library the singular-value decomposition routine DGESVD() would be used.} is then used to compute the null-space of $A$, which is placed in a $4\times 4$ matrix $B$. The signs of the entries in the first column of $B$ are  flipped and the resulting matrix $B$ is then used to operate, by standard matrix multiplication, on all the quaternions generated by rejection sampling to move the cluster of quaternions so that $\bm{\mu}$ becomes the cluster's mean direction.

As an example, we use the rejection sampling algorithm described above to construct two samples of orientations:
\begin{itemize}
	\item A random sample of $1,000$ orientations clustered around the $120^{\circ}@[111]$ orientation with $\Delta\theta=10^{\circ}$, using the vMF pdf;
	\item A random sample of $100$ orientations clustered around the $90^{\circ}@[0\bar{1}0]$ orientation with $\Delta\theta=5^{\circ}$, using the Watson pdf.
\end{itemize}
The corresponding $\kappa$ value is equal to $\kappa=1/(1-\cos\Delta\theta)$, which is $65.8$ for the vMF case, and $262.8$ for the Watson case.  The means directions are represented by the quaternions 
\[
	\mu_{\text{vMF}}=\left[\frac{1}{2},\frac{1}{2},\frac{1}{2},\frac{1}{2}\right]
\]
and 
\[
	\mu_{\text{Watson}}=\left[\frac{1}{\sqrt{2}},0,\frac{1}{\sqrt{2}},0\right]\ .
\]
The two distributions are shown as 3D stereographic projections in Fig.~\ref{\chabbr:fig:samplesSP}, the vMF case in (a), and the Watson case in (b).  

\begin{figure}[h!]
\centering\leavevmode
\includegraphics[width=0.9\textwidth]{\graphicspath samplesSP.pdf}
\caption{\label{\chabbr:fig:samplesSP}(a) Stereographic projection of the von Mises-Fisher random sampling described in the text; (b) same for the Watson random sampling.}
\end{figure}

\subsection{Orientation averaging on $\mathbb{S}^3$}\label{\chabbr:sec:oraverages3}
In the previous section we answered the question: given a mean orientation and a concentration parameter, generate a random orientation sample using either the vMF or Watson pdf.  In this section we turn the question around: given a set of orientations, what are the mean orientation and concentration parameter that best describe this set for either the vMF or the Watson pdf?  It is important to note that this section deals with the average orientation \textit{in the absence of crystallographic symmetry}; the averaging process becomes significantly more involved when fundamental zones are considered.

\subsubsection{The quaternion logarithm average}\label{\chabbr:ssec:quatlogav}
We have seen that, in the context of orientations, eq.~(\ref{\chabbr:eq:regularaverage}) does not necessarily provide the correct answer for the average orientation due to the double-cover nature of the unit quaternion sphere $\mathbb{S}^3$. In particular, if we double the size of the set $\{q_i\}$ by adding in the $\{-q_i\}$ rotations, then the simple arithmetic mean of eq.~(\ref{\chabbr:eq:regularaverage}) will produce a mean resultant length of $\bar{R}=0$, and the mean unit quaternion $\bm{\mu}$ is undefined.  Instead of the conventional mean defined in eq.~(\ref{\chabbr:eq:regularaverage}), we can use the quaternion logarithm and exponential operations to obtain an average. Euler's formula for 2D rotations, $\mathrm{e}^{\mathrm{i}\theta}=\cos\theta+\mathrm{i}\sin\theta$, has a quaternion equivalent of the form:
\begin{equation}
	\mathrm{e}^{(n_x\mathrm{i}+n_y\mathrm{j}+n_z\mathrm{k})\frac{\theta}{2}} = \cos\frac{\theta}{2}+(n_x\mathrm{i}+n_y\mathrm{j}+n_z\mathrm{k})\sin\frac{\theta}{2}\ .
\end{equation}
Applying logarithm rules to the exponential on the left results in:
\begin{equation}
	\ln\left(\mathrm{e}^{(n_x\mathrm{i}+n_y\mathrm{j}+n_z\mathrm{k})\frac{\theta}{2}}\right) = (n_x\mathrm{i}+n_y\mathrm{j}+n_z\mathrm{k})\frac{\theta}{2}\equiv
	\hat{\mathbf{n}}_j\frac{\theta_j}{2}\ ;
\end{equation}
note that the right-hand side, for some orientation $j$ in a set of orientations, is similar to the axis-angle representation, but with half the rotation angle. This is a pure quaternion with zero scalar part.  We can then compute the conventional geometric mean of the right-hand side over all orientations in the set, and then take the exponential of the result to obtain an estimate of the average quaternion. If we set:
\[
	\mathbf{s} = \frac{1}{n}\sum_{j=1}^{n} \hat{\mathbf{n}}_j\frac{\theta_j}{2}
\]
then we can take the quaternion exponential by computing the length of $\mathbf{s}$ as $\bar{s}\equiv\vert\mathbf{s}\vert$, and using:
\[
	\bm{\mu} = \left[ \cos(\bar{s}), [s_1,s_2,s_3]\frac{\sin(\bar{s})}{\bar{s}} \right]\ .
\]
The average rotation angle, $\omega$, is then $\omega=2\bar{s}$. 

Once again we can illustrate the need to work only with quaternions that belong to one hemisphere of $\mathbb{S}^3$.  If we just have two quaternions, $q$ and $-q$, then the standard geometric average results in $\frac{1}{2}(q-q)=0$, i.e., the zero quaternion, which does not represent a rotation. If instead we use the logarithmic average, we have:
\[
	\ln(q) = \hat{\mathbf{n}}\frac{\theta}{2}\text{ and } \ln(-q) = -\hat{\mathbf{n}}\frac{\theta}{2}\rightarrow \mathbf{s} = \frac{1}{2}(\ln(q)+\ln(-q)) = (0,0,0)\ ;
\]
in that case we have $\bar{s} = 0$, so that $\cos\bar{s}=1$ and $\frac{\sin\bar{s}}{\bar{s}}=1$ (in the limit of $\bar{s}\rightarrow 0$), and the resulting quaternion exponential is $\bm{\mu}=(1,0,0,0)$ which is the identity quaternion.  While the result of the logarithmic averaging of $q$ and $-q$ is a unit quaternion, which is not always the case for the standard geometric mean, the fact that the logarithm average is the identity quaternion for every antipodal pair illustrates that this procedure does not produce any meaningful results if the full $\mathbb{S}^3$ hypersphere is allowed. 

\subsubsection{The scatter matrix average}\label{\chabbr:ssec:scatav}
We have introduced the scatter matrix before as the sum of the outer products of all the quaternions in the set $\{q_j\}$. We can slightly generalize this approach by allowing weight factors $w_j$ for the individual quaternions in the set (with the restriction that all quaternions belong to $\mathbb{S}^3_+$ and in the absence of crystallographic symmetry).  Generating the scatter matrix as before we obtain:
\begin{equation}
	T_{kl} = \sum_{j=1}^n w_j q_{j,k}q_{j,l}\ ,
\end{equation}
where $q_{j,k}$ is the $k$-th component of quaternion $q_j$.  One can then show that the eigenvector corresponding to the largest eigenvalue of $T_{kl}$ is the weighted mean direction $\bm{\mu}$ of the set $\{q_j\}$.  

As a simple example, consider two $90^{\circ}$ rotations, one around the $[100]$ axis, the other around $[001]$. What is the average rotation?  Putting $s=1/\sqrt{2}$, the quaternions are given by $q_1=[s,s,0,0]$ and $q_2=[s,0,0,s]$. Taking unit weights, the scatter matrix is then readily computed to be:
\[
	T = q_1\otimes q_1+q_2\otimes q_2 = \begin{pmatrix} 1&1/2&0&1/2\\ 1/2&1/2&0&0\\ 0&0& 0&0\\ 1/2&0&0&1/2\end{pmatrix}
\]
The eigenvalues of this matrix are (in descending order) $[3/2,1/2,0,0]$ and the eigenvectors are given by the columns in the following matrix (same order as the eigenvalues):
\[
	\begin{pmatrix}2&0&-1&0\\ 1&-1&1&0\\ 0&0&0&1\\1&1&1&0\end{pmatrix}\ .
\]
The first column, when normalized, results in the quaternion 
\[
	q_{\text{av}}=\frac{1}{\sqrt{6}}[2,1,0,1]=[0.816497, 0.408248, 0.000000, 0.408248]\ ,
\]
which represents a rotation of $70.53^{\circ}$ around the $[101]$ axis, i.e, the bisector of the individual rotation axes. This is in full agreement with the standard geometric mean $q_{\text{av}}=\frac{1}{2}(q_1+q_2)$. 


\subsubsection{von Mises-Fisher and Watson averaging}\label{\chabbr:ssec:vMFWav}

We start with a probability density function $f(x;\theta)$, where $x$ is a realization of a random variable $X$, and $\theta$ is a set of one or more parameters that the pdf depends on.  In the usual interpretation of a pdf, we think of the function as being evaluated for a fixed set of parameters $\theta$, and $x$ is varied by taking different samples from $X$. On the other hand, if we fix the sample $x$, and we let the parameter $\theta$ vary, then $f(x;\theta)$ is interpreted as the \indexit{likelihood function}, and by maximizing the likelihood with respect to the parameters, we can determine the value of $\theta$ that best matches the pdf for the selected sample $x$.

Let us illustrate this concept by means of a simple 1-D example. Consider a set of $N=10$ temperature measurements, $T_i$, inside an annealing furnace using a thermocouple that has a stated accuracy of $\pm 0.1$K. The measurements are carried out at random times over a $24$-hour period, and the measurement conditions are identical for each measurement. The set-point of the furnace is the temperature $T_s=773.1$K (or $500^{\circ}$C), and the distribution of the measured values is expected to be normal with a narrow variance, $\sigma_T$.  A representative set of measurements is shown in Table~\ref{\chabbr:tb:thermal}.

% table
\begin{table}[ht]
\centering\leavevmode
\begin{tabular}{|l|cccccccccc|}
\hline
$i$ & 1 & 2 & 3 & 4 & 5 & 6 & 7 & 8 & 9 & 10 \\
$T_i$ (K) &773.6   &   776.6 &     770.6 &      774.5 &      778.9 &      776.6 &      774.8 &      774.3 &      785.4 &      766.4\\
\hline
\end{tabular}
\caption{\label{\chabbr:tb:thermal}Example set of temperature measurements.}
\end{table}

Instead of maximizing the likelihood function with respect to the parameters, it is common practice to maximize the \indexit{log-likelihood} function to remove the exponential function; this is obviously permitted since the logarithm is a monotonically increasing function. If we start with a standard 1-D normal distribution:
\[
	f(T;T_\mu,\sigma_T^2) = \frac{1}{\sqrt{2\pi\sigma_T^2}}\mathrm{e}^{-\frac{(T-T_\mu)^2}{2\sigma_T^2}}\ ,
\]
then the likelihood function for a sample $T=\{T_1,\ldots,T_N\}$ of  independent and identically distributed (i.i.d.) measurements is given by:
\[
	\mathcal{L}(T_\mu,\sigma_T^2) =  \prod_{i=1}^N f(T_i;T_\mu,\sigma_T^2)\ .
\]
Therefore, the log-likelihood is given by (using standard logarithm rules):
\begin{equation*}
\begin{split}
	\ln \mathcal{L}(T_\mu,\sigma_T^2) &= \sum_{i=1}^N \ln  f(T_i;T_\mu,\sigma_T^2)\ ;\\
	&= -\frac{N}{2}\ln(2\pi) - \frac{N}{2}\ln(\sigma_T^2) -\frac{1}{2\sigma_T^2}  \sum_{i=1}^N (T_i-T_{\mu})^2\ .
\end{split}
\end{equation*} 
The quantities $(T_i-T_{\mu})$ are known as the \indexit{residuals}.  The parameters that maximize the log-likelihood are then obtained by setting the partial derivatives with respect to each parameter equal to zero. We find for $T_{\mu}$:
\[
	\frac{\partial \ln \mathcal{L}(T_\mu,\sigma_T^2)}{\partial T_{\mu}} = \frac{1}{\sigma_T^2}  \sum_{i=1}^N (T_i-T_{\mu}) = 0\ ,
\]
(i.e., the sum of the residuals is zero) from which we derive:
\begin{equation}
	T_{\mu} = \frac{1}{N} \sum_{i=1}^N T_i\ ,
\end{equation}
which is the standard arithmetic mean of the observations. For the example in Table~\ref{\chabbr:tb:thermal} we have $T_{\mu} = 775.2$ K.
For the standard deviation, $\sigma_T^2$, we have:
\[
	\frac{\partial \ln \mathcal{L}(T_\mu,\sigma_T^2)}{\partial \sigma_T^2} =  - \frac{N}{2\sigma_T^2} +\frac{1}{2\sigma_T^4}  \sum_{i=1}^N (T_i-T_{\mu})^2 = 0\ ,
\]
from which we derive an expression for the standard deviation $\sigma_T^2$:
\begin{equation}
	\sigma_T^2 = \frac{1}{N} \sum_{i=1}^N (T_i-T_{\mu})^2\ .
\end{equation}
For the example data set, we find $\sigma_T^2=22.43$ K$^2$ or $\sigma_T=4.7$ K.  

For small sample sizes, the expression for the standard deviation is often multiplied by the correction factor $N/(N-1)$, which accounts for the fact that, while the measurements $T_i$ are independent, the residuals $(T_i-T_{\mu})$ are not independent since they sum to zero.  Incorporating the correction into the expression for the standard deviation results in the \indexit{unbiased estimator} $s_T$:
\begin{equation}
	s_T^2 = \frac{1}{N-1} \sum_{i=1}^N (T_i-T_{\mu})^2\ ;
\end{equation}
for the example data we find $s_T^2= 24.9$ K$^2$ and $s_T=5.0$ K.  For more details on unbiased estimators we refer the interested reader to standard textbooks on statistics.

The same log-likelihood procedure as described above can be carried out for the vMF and Watson pdfs to obtain expressions for the mean quaternion, $\bm{\mu}$, and the concentration parameter, $\kappa$.  If all the observed orientations, $\{q_i\}$, are clustered into a single cluster, then this log-likelihood approach will produce the proper parameters $\bm{\mu}$ and $\kappa$.  If, on the other hand, there are multiple clusters, then it may not always be clear a-priori to which cluster a given orientation quaternion belongs, and the problem becomes more involved.  Let us first address the problem of a single cluster with $N$ orientations.

The vMF likelihood function for a sample $\{q_1,q_2,\ldots,q_N\}$ of $N$ independent and identically distributed (i.i.d.) orientation quaternions is given by:
\[
	\mathcal{L}(\bm{\mu},\kappa) = \prod_{i=1}^N \frac{\kappa}{(2\pi)^2I_1(\kappa)}\mathrm{e}^{\kappa\,\bm{\mu}\cdot q_i} = \left( \frac{\kappa}{(2\pi)^2I_1(\kappa)}\right)^N
	\mathrm{e}^{\kappa\,\bm{\mu}\cdot \sum_{i=1}^N q_i} \ .
\]
In this relation, we consider the quaternions in the exponential as vectors in $\mathbb{R}^4$.  The sum in the exponent, by analogy with the analysis in section~\ref{\chabbr:ssec:circleaveraging}, is the \indexit{resultant vector} of the input quaternions, $\mathbf{Q}=\sum_{i=1}^N q_i$. The vMF likelihood function thus becomes:
\[
	\mathcal{L}(\bm{\mu},\kappa) = \left(\frac{\kappa}{(2\pi)^2I_1(\kappa)}\right)^N \mathrm{e}^{\kappa\,\bm{\mu}\cdot\mathbf{Q}} \ .
\]
The log-likelihood function is then readily shown to be:
\[
	\ln\left(\mathcal{L}(\bm{\mu},\kappa)\right) = N[\ln\kappa -2\ln(2\pi)- \ln(I_1(\kappa))] + \kappa\, \bm{\mu}\cdot \mathbf{Q}\ .
\]
Next, we maximize this function with respect to $\bm{\mu}$ and $\kappa$. The parameter $\bm{\mu}$ only appears in the last term, so we need to maximize this term, subject to the constraint that $\vert \bm{\mu}\vert=1$.  The Cauchy-Schwarz inequality states:
\[
	\vert \bm{\mu}\cdot \mathbf{Q}\vert \le \vert\bm{\mu}\vert\,\vert\mathbf{Q}\vert = Q\ ,
\]
where $Q=\vert\mathbf{Q}\vert$. Equality only occurs when $\bm{\mu}= \mathbf{Q}/Q$, so that the estimate for the mean vMF orientation is given by:
\[
 	\hat{\bm{\mu}} = \mathbf{Q}/Q = \frac{\sum_{i=1}^N q_i}{\vert \sum_{i=1}^N q_i\vert}\ .
\]
Note that we can substitute this estimate into the log-likelihood function to obtain:
\[
	\ln\left(\mathcal{L}(\hat{\bm{\mu}},\kappa)\right) = N[\ln\kappa -2\ln(2\pi)- \ln(I_1(\kappa))] + \kappa\, Q\ .
\]
The second parameter, $\kappa$, occurs several times in the log-likelihood function, so we take the partial derivative with respect to $\kappa$ and set it equal to zero:
\[
	\frac{\partial\ln\left(\mathcal{L}(\hat{\bm{\mu}},\kappa)\right)}{\partial\kappa} = N\left[\frac{1}{\kappa}-\frac{I'_1(\kappa)}{I_1(\kappa)}\right]+Q\ .
\]
Using a standard identity for modified Bessel functions we can eliminate the derivative $I'_1(\kappa)$ and obtain:
\[
	\frac{\partial\ln\left(\mathcal{L}(\hat{\bm{\mu}},\kappa)\right)}{\partial\kappa} = N\left[\frac{1}{\kappa}-\left(\frac{I_2(\kappa)}{I_1(\kappa)}+\frac{1}{\kappa}\right)\right]+Q\ .
\]
Simplifying, and putting the derivative equal to zero, we obtain a transcendental equation for the estimated $\hat{\kappa}$:
\[
	\frac{I_2(\hat{\kappa})}{I_1(\hat{\kappa})} = \frac{Q}{N}\ .
\]
This equation must be solved numerically for the estimated concentration parameter $\hat{\kappa}$. The maximum likelihood estimate for the vMF distribution, given a set of i.i.d.\ orientations $\{q_i\}$, is thus given by:
\[
	f(q;\hat{\bm{\mu}},\hat{\kappa}) = \frac{\hat{\kappa}}{(2\pi)^2I_1(\hat{\kappa})}\mathrm{e}^{\hat{\kappa}\,\hat{\bm{\mu}}\cdot q}\ .
\]

A similar process can be followed for the Watson distribution, but the derivation is slightly more complex due to the quadratic nature of the argument of the exponential (eq.~\ref{\chabbr:eq:Watsonpdf}): 
\[
	f_W(q;\bm{\mu},\kappa) = \frac{\mathrm{e}^{\kappa(\bm{\mu}\cdot q)^2}}{{_1}F_1[\frac{1}{2},2,\kappa]}=\frac{\mathrm{e}^{-\frac{\kappa}{2}}}{I_0(\frac{\kappa}{2})-I_1(\frac{\kappa}{2})}\mathrm{e}^{\kappa(\bm{\mu}\cdot q)^2}
\]
The denominator in the first equality is the \indexit{Kummer confluent hypergeometric function}, which can be expanded in terms of modified Bessel functions.  For an i.i.d.\ set of $N$ orientations $\{q_i\}$, the likelihood function is given by:
\[
	\mathcal{L}(\bm{\mu},\kappa) = \prod_{i=1}^N f_W(q;\bm{\mu},\kappa) = \left(\frac{1}{{_1}F_1[\frac{1}{2},2,\kappa]}\right)^N \mathrm{e}^{\kappa\sum_{i=1}^N(\bm{\mu}\cdot q_i)^2}\ ,
\] 
which results in the following log-likelihood function:
\[
	\ln(\mathcal{L}(\bm{\mu},\kappa)) = -N\ln\left[ {_1}F_1[\frac{1}{2},2,\kappa] \right] + \kappa\sum_{i=1}^N(\bm{\mu}\cdot q_i)^2\ .
\] 
Since we have
\[
	\sum_{i=1}^N(\bm{\mu}\cdot q_i)^2 = \sum_{i=1}^N (\bm{\mu}\cdot q_i)\otimes(q_i \cdot \bm{\mu}) = \bm{\mu}\cdot \left(\sum_{i=1}^N q_i\otimes q_i\right) \cdot \bm{\mu}\ ,
\]
we see that the right-hand side involves the scatter matrix, $\mathbf{T}$, introduced in section~\ref{\chabbr:ssec:scatav}:
\[
	\sum_{i=1}^N(\bm{\mu}\cdot q_i)^2 = \bm{\mu}^T\mathbf{T}\bm{\mu}\ ,
\]
where the transpose symbol indicates the vector is written as a row instead of a column. The log-likelihood function can thus be written as:
\[
	\ln(\mathcal{L}(\bm{\mu},\kappa)) = -N\ln\left[ {_1}F_1[\frac{1}{2},2,\kappa] \right] + \kappa\,\bm{\mu}^T\mathbf{T}\bm{\mu}\ .
\] 
As for the vMF case, the mean direction, $\bm{\mu}$, is constrained to be a unit quaternion, $\vert\bm{\mu}\vert=1$.  

Next we must maximize the log-likelihood function with respect to the two parameters.  The parameter $\bm{\mu}$ only occurs in the last term, so we must maximize this term. Since there is a constraint on the magnitude of $\bm{\mu}$, we can use the method of \indexit{Lagrange multipliers}:
\[
	L = \kappa\,\bm{\mu}^T\mathbf{T}\bm{\mu} - \lambda(\bm{\mu}^T\bm{\mu}-1)\ .
\]
Setting the derivative with respect to $\bm{\mu}$ equal to zero we have
\[
	\frac{\mathrm{d}L}{\mathrm{d}\bm{\mu}} = 2\kappa \mathbf{T}\bm{\mu}-2\lambda \bm{\mu} = 0\ ,
\]
where the factors of $2$ stem from the quadratic nature of both terms.  We find that:
\[
	\mathbf{T}\bm{\mu} = \frac{\lambda}{\kappa}\bm{\mu}\ ,
\]
which indicates that $\bm{\mu}$ must be an eigenvector of the scatter matrix $\mathbf{T}$.  Which of the four eigenvectors we select depends on the sign of $\kappa$; if $\kappa>0$, then we select the eigenvector corresponding to the largest eigenvalue, and we have an axial distribution. If $\kappa<0$, then we select the eigenvector with the smallest eigenvalue, and we have a so-called \indexit{girdle distribution}. The mean direction estimate $\hat{\bm{\mu}}$ is therefore either the principal eigenvector of $\mathbf{T}$ or the minor eigenvector, depending on the sign of $\kappa$. We introduce the parameter $\Lambda$, defined as:
\[
	\Lambda = \hat{\bm{\mu}}^T\mathbf{T}\hat{\bm{\mu}}\ ;
\]
the log-likelihood function now becomes:
\[
	\ln(\mathcal{L}(\hat{\bm{\mu}},\kappa)) = -N\ln\left[ {_1}F_1[\frac{1}{2},2,\kappa] \right] + \kappa\Lambda\ ,
\] 
and we need to maximize this function with respect to the concentration parameter, $\kappa$, to complete the MLE computation.  The derivative with respect to $\kappa$ is given by:
\[
	\frac{\partial \ln(\mathcal{L}(\hat{\bm{\mu}},\kappa)) }{\partial\kappa} = -N\frac{ {_1}F'_1[\frac{1}{2},2,\kappa] }{ {_1}F_1[\frac{1}{2},2,\kappa] }+\Lambda = 0\ ,
\]
so that:
\[
	\frac{ {_1}F'_1[\frac{1}{2},2,\kappa] }{ {_1}F_1[\frac{1}{2},2,\kappa] } = \frac{\Lambda}{N}\ .
\]
This is another transcendental equation that must be solved numerically for $\kappa$.  The left-hand side of this equation can be expressed in terms of modified Bessel functions and we obtain:
\[
	\frac{I_1(\frac{\hat{\kappa}}{2})}{I_1(\frac{\hat{\kappa}}{2})+I_0(\frac{\hat{\kappa}}{2})} = -\frac{\hat{\kappa}\Lambda}{N}\ .
\]
This completes the MLE treatment for the Watson distribution.  We should re-emphasize that the above derivations are only valid for orientation data sets in the absence of crystallographic symmetry; when rotational symmetry, and thus fundamental zones, are included, the statistical averaging procedure becomes more involved, as will be described in the next section.



\section{Orientation Sampling and Averaging in the Presence of Fundamental Zones}\label{\chabbr:sec:samplingRFZs}
In section~\ref{\chabbr:ssec:randomsampling}, several approaches are described to generate random samples of orientations; Schoemake sampling and Super-Fibonacci sampling produce high quality (i.e., uniform) orientation sets, whereas others, like Marsaglia sampling, do not have perfect uniformity, but are still reasonably uniform.  A new orientation representation, the \indexit{cubochoric representation} is introduced in section~\ref{3DROTS:sec:cubochoric}; this representation was developed with the intent of generating uniform samples of orientations in a numerically convenient way, by starting from a uniform cubical grid.  In the present section, we answer two questions: how are these orientation sampling approaches affected by the introduction of Fundamental Zones? And, how do we average orientations in the presence of crystallographic symmetry?

\subsection{Orientation sampling and crystallographic symmetry}\label{\chabbr:ssec:RFZsampling}
When the cubochoric representation is combined with the Rodrigues representation, one can generate uniform orientation samples that are restricted to a crystallographic Fundamental Zone, i.e., perform \indexit{rejection sampling} for which being inside or outside of the relevant FZ is the deciding factor.  Consider the following two examples: the point group symmetries are \pg{24} and \pg{30}, with $N=20$ in the first case and $N=50$ in the second. Without the application of symmetry the total number of points generated would be $(2N+1)^3$ or $68,921$ for $N=20$ and $1,030,301$ for $N=50$.  The application of symmetry is performed by a simple test: first convert each cubochoric point to a Rodrigues vector, then test whether or not the point lies inside the RFZ; keep the point if it lies inside the RFZ, discard it otherwise. The resulting lists of orientations have $n=5,367$ for \pg{24} and $n=41,625$ for \pg{30}.  The distributions of points along with the RFZ edges are shown as stereographic projections in Fig.~\ref{\chabbr:fig:samples}. Note in particular for \pg{30} that the sampling points lie on concentric spherical surfaces inside the RFZ due to the way they are sampled on concentric cubes in the cubochoric cell.

\begin{figure}[h!]
\centering\leavevmode
\includegraphics[width=0.9\textwidth]{\graphicspath samples.pdf}
\caption{\label{\chabbr:fig:samples}Stereographic projection of orientations uniformly sampled in cubochoric space and restricted to the Rodrigues Fundamental Zone for \pg{24} and \pg{30}.}
\end{figure}

When this type of orientation sampling is subjected to tests of uniformity, for instance through comparison of the Riesz energies with the theoretical optimal values, or an analysis of the dispersion of the set of orientations, it is found that, near the boundaries of the fundamental zone, the local uniformity depends strongly on the choice of the number $N$. With increasing $N$, the effect of this local non-uniformity decreases significantly but it never completely disappears.  Resolving this issue requires a symmetry-aware remapping of the orientation set using an approach that involves  Knothe--Rosenblatt (KR) rearrangements in spherical homochoric coordinates.% \cite{xxx}.

\subsection{Orientation averaging in the presence of crystallographic symmetry}\label{\chabbr:ssec:RFZaveraging}

Before we describe von Mises-Fisher and Watson averaging in the presence of rotational symmetry, let us first illustrate by means of an example why fundamental zones complicate the process.  We begin by creating an orientation sample with $1,000$ orientations around the rotation $45^{\circ}@[010]$ in a cubic (octahedral) material.  We generate the sample using the von Mises-Fisher sampling approach for a concentration parameter value of $\kappa=250$ (corresponding to a misorientation angle spread of about $5^{\circ}$) without application of the fundamental zone. The resulting cluster is shown in the stereographic representation in Fig.~\ref{\chabbr:fig:Goss}(a).  Application of the scatter matrix averaging procedure to this data set results in a mean rotation quaternion of $\hat{\bm{\mu}}=[0.9243245, 0.0001555,  0.3815980,-0.0026589]$, which is a rotation by $\omega=44.87^{\circ}$ around the axis $[0.00041,0.99997,-0.00696]$, different from $[010]$ by only $0.399^{\circ}$.  

\begin{figure}[h!]
\centering\leavevmode
\includegraphics[width=0.75\textwidth]{\graphicspath Goss.pdf}
\caption{\label{\chabbr:fig:Goss}Stereographic projection of $1,000$ orientations sampled using a von Mises-Fisher distribution with $\bm{\mu}=[0,1,0]$ and $\kappa=250$; in (a) no symmetry is applied whereas (b) shows a split cluster due to the octahedral fundamental zone.}
\end{figure}

Next we apply the octahedral fundamental zone and we find that the orientation cluster is now split into two parts by the RFZ boundary (Fig.~\ref{\chabbr:fig:Goss}(b)). If we once again use the scatter matrix approach to average the orientations we find that the mean rotation is very close to $\hat{\bm{\mu}}=[1,0,0,0]$, which is the identity rotation. This makes sense since the two halves of the original cluster are now symmetrically located in opposite parts of the fundamental zone so that their average is a point near the center of the RFZ.  This example clearly shows that averaging in the presence of rotational crystallographic symmetry will require a more advanced method.  

\begin{figure}[h!]
\centering\leavevmode
\includegraphics[width=0.75\textwidth]{\graphicspath Goss2.pdf}
\caption{\label{\chabbr:fig:Goss2}Stereographic projection of $1,000$ orientations sampled using a von Mises-Fisher distribution with $\bm{\mu}=[0,1,0]$ and $\kappa=250$; in (a) octahedral symmetry is applied whereas (b) shows a similar distribution with $\kappa=5,000$.}
\end{figure}

We will work primarily with the von Mises-Fisher distribution, but an analogous procedure can be applied to the Watson distribution.  The example in Fig.~\ref{\chabbr:fig:Goss} shows that simple application of the FZ renders it impossible to use the standard averaging procedures, including the MLE vMF/Watson approaches, because there can be more than one clusters intersecting the FZ boundaries. We can remedy this problem by  applying the symmetry operators to the orientation dataset.  Fig.~\ref{\chabbr:fig:Goss2}(a) shows the symmetrized version of the orientation data set used in Fig.~\ref{\chabbr:fig:Goss}; to more clearly show the locations of the $24$ clusters, Fig.~\ref{\chabbr:fig:Goss2}(b) shows a dataset with a concentration parameter of $\kappa=5,000$, which yields much tighter clusters.  

The solution to estimating the mean orientation $\hat{\bm{\mu}}$ and concentration parameter $\hat{\kappa}$ is to make the probability density function itself rotationally invariant with respect to the point group symmetry of the crystal structure.  The invariant von Mises-Fisher pdf is then given by a mixture of individual vMFs, one for each symmetry operator:
\begin{equation}
	\tilde{f}(q;\bm{\mu},\kappa) = \frac{1}{M} \sum_{m=1}^M \frac{\kappa}{4\pi^2 I_1(\kappa)}\mathrm{e}^{\kappa\,\bm{\mu}\cdot (Q_m q)}\ ,\label{\chabbr:eq:symvMFpdf}
\end{equation}
where $M$ is the order of the rotational point group with symmetry operators $Q_m$. In this form of the invariant pdf, the experimental orientations $\{q_i\}$ are symmetrized into multiple equivalent clusters, as shown in Fig.~\ref{\chabbr:fig:Goss2}(a). This is entirely equivalent to applying the operators $Q_m$ to the mean direction $\bm{\mu}$: $\bm{\mu}_m = Q_m\bm{\mu}$; in other words, the mean directions of the clusters are related to one another by the symmetry operators:
\begin{equation}
	\tilde{f}(q;\bm{\mu},\kappa) = \frac{1}{M} \sum_{m=1}^M \frac{\kappa}{4\pi^2 I_1(\kappa)}\mathrm{e}^{\kappa\bm{\mu}_m\cdot q}\ ,\label{\chabbr:eq:symvMFpdf2}
\end{equation}
Since this pdf has $M$ clusters, i.e., it is a \textit{mixture} of $M$ individual vMF pdfs, there is no way of knowing a-priori which of these clusters the input data $\{q_i\}$ belong to; there is therefore missing information in this problem, and we typically refer to such data as \indexit{hidden data} or \indexit{latent variables}.  

The determination of the pdf parameters in the presence of hidden data is typically approached by means of the \indexit{Expectation Maximization} or EM algorithm \cite{dempster1977a}.  This is a two-step iterative algorithm in which one alternates between estimating the probability of each of the hidden variables (the E-step, which guesses which cluster each data point belongs to, given the current model parameters) and estimating the parameters (using the probabilities from the E-step) to maximize the likelihood of the data (this is the M-step, which maximizes the fit).  

Next, we proceed as before by writing down the likelihood and log-likelihood functions (we represent the pre-factor of the vMF pdf by the symbol $c_4(\kappa)$):
\[
	\mathcal{L}(\bm{\mu},\kappa)) = \prod_{i=1}^N \sum_{m=1}^M c_4(\kappa)\mathrm{e}^{\kappa\bm{\mu}_m\cdot q_i}\ ;
\]
\[
	\ln\left(\mathcal{L}(\bm{\mu},\kappa))\right) = \sum_{i=1}^N \sum_{m=1}^M \left(\ln(c_4(\kappa))+\kappa\bm{\mu}_m\cdot q_i\right)\ .
\]
Since we wish to maximize the likelihood, we have omitted the constant factor $1/M$ from the expressions.  In the EM algorithm, the function to be maximized is a weighted form of the log-likelihood function to account for the latent parameters:
\[
	Q(\bm{\mu},\kappa)) = \sum_{i=1}^N \sum_{m=1}^M r_{i,m}\,\left(\ln(c_4(\kappa))+\kappa\bm{\mu}_m\cdot q_i\right)\ .
\]
The factors $r_{i,m}$ represent the (posterior) probability that the orientation $q_i$ belongs to cluster $m$, given the parameters $(\bm{\mu}_m,\kappa)$.  In the E-step, we set:
\[
	r_{i,m} = \frac{f(q_i;\bm{\mu}_m,\kappa)}{\sum_{l=1}^M f(q_i;\bm{\mu}_l,\kappa)}\,\text{ where }m\in\{1,2,3,\ldots,M\}\ .
\]
In the M-step, we apply the standard maximum likelihood estimation to the function $Q$, by computing the partial derivatives with respect to both parameters and setting them equal to zero. This derivation is identical to the one in section~\ref{\chabbr:ssec:vMFWav}, and the resulting EM estimate can be expressed by defining the resultant quaternion as:
\[
	\bm{\gamma} \equiv \sum_{i=1}^N \sum_{m=1}^M r_{i,m}\, Q_m^T q_i\ ,
\]
so that the mean direction is given by:
\[
	\hat{\bm{\mu}} = \frac{\bm{\gamma}}{\vert\bm{\gamma}\vert}\ ,
\]
and $\hat{\kappa}$ can be determined numerically by solving the transcendental equation:
\[
	\frac{I_2(\hat{\kappa})}{I_1(\hat{\kappa})} = \frac{N\hat{\kappa}}{\vert\bm{\gamma}\vert}\ .
\]
This process is iterated until the changes in the estimated parameters $(\hat{\bm{\mu}},\hat{\kappa})$ become smaller than some preset threshold. For a mixture of Watson distributions, a similar approach is followed to determine the parameter estimates; details can be found in \cite{degraef2015zg}.

Application of the mixture of vMF pdfs EM approach to the data set used in Fig.~\ref{\chabbr:fig:Goss}(b), and employing the octahedral symmetry group, we find for the estimated parameters:
\[
	\hat{\bm{\mu}} = [0.92431670,-0.00137444, 0.38162359,-0.00041897]\text{ and }\hat{\kappa}=251.65\ ;
\]
This quaternion has a rotation angle that is only $0.21^{\circ}$ different from $45^{\circ}$, and the rotation axis is also $0.21^{\circ}$ different from $[010]$.  The quaternion logarithm and the scatter matrix algorithms both produce the identity rotation, so those approaches are clearly not able to handle the presence of rotational symmetry.



